\clearpage\begingroup\SourceLeaf{88b}\HistoricalHead{83}{R}{GAUSS AN ENCKE}\HistoricalFoot{}{11\textsuperscript{*}}
\begin{center}[8.]\end{center}
\begin{center}\textsc{Gauss} an \textsc{Encke}.\quad Göttingen, 18. August 1832.\end{center}
\noindent . . . . . . . . . . . . . . . . . . . . .\par
\indent Meine Beschäftigungen mit dem Magnetismus haben seit meinem letzten
Briefe fortgedauert; meine Apparate, die ich in Duplo fertigen zu lassen für
nöthig gehalten habe, sind in sehr vielen Stücken weiter vervollkommnet, und
es bleibt jetzt eigentlich gar nichts weiter zu wünschen übrig, als ein gegen
Eisennähe und Luftzug ganz geschütztes Local. Mein Hauptaugenmerk ist
noch auf die absolute Intensität gerichtet, und ich habe die Absicht, über
diesen Theil der Untersuchungen zuerst etwas öffentlich zu sagen, wahrscheinlich
in einer Societätsvorlesung. Binnen einigen Monaten wird sich vielleicht
die Ausarbeitung, mit der ich bereits einen kleinen Anfang gemacht habe,
vollenden lassen. Bestimmt kann ich aber darüber nichts sagen. Ich hasse
alles übereilte Publiciren und wünsche immer nur reifes zu geben, und da
trifft es sich denn nicht selten, dass wegen dieses oder jenen Umstandes, der,
\textls[120]{nachdem} er erledigt ist, wenige Zeilen füllt, ein wochen- oder monatelanger
Aufenthalt entsteht. In dem bisherigen Vortrag der Lehre vom Magnetismus
findet sich aber so viel Vages, Nichtssagendes, Unlogisches (auch selbst bei
\textsc{Biot}[\textsuperscript{*})]),\begingroup\renewcommand{\thefootnote}{}
\footnotetext{[\textsuperscript{*}) Gemeint ist wohl der dritte Band von J. B. \textsc{Biot}, \textit{Traité de Physique expérimentale et mathématique}, 4 vol, Paris 1816.]}
\endgroup\ dass hier erst ganz von vorne an aufgebaut werden muss. Es gehört
dahin der \textls[120]{Begriff} der Pole. Dann der schreiende Widerspruch, dass man
einmal annimmt, \textls[120]{in jedem Theilchen} einer Nadel sei eben so viel nördlicher
als südlicher Magnetismus, und nachher doch immer so spricht als sei
an einem Ende der Nadel bloss der Eine am andern der andere Magnetismus.
Mich hat diese Verworrenheit bei \textsc{Biot} im vorigen Herbst, als ich anfing mich
mit diesen Dingen zu beschäftigen, erst lange gequält. Ich konnte mit \textls[120]{seinem}
freien Magnetismus gar keinen Sinn verbinden. Durch die Beziehung auf die
Electricität hat \textsc{Biot} die Sache nur verwirrter gemacht. Ich bin nun freilich
\clearpage\SourceLeaf{89}\HistoricalHead{84}{L}{PHYSIK. BRIEFWECHSEL.}
\noindent in diesen Dingen schon lange zu völliger Klarheit gekommen, allein es gehören
dazu mehrere neue höchst interessante Lehrsätze die \textls[120]{sehr tief} liegen, deren
\textls[120]{Entwicklung} die Grenzen der mir zunächst vorgesetzten Abhandlung weit,
sehr weit überschreiten würden, und die ich daher \textls[120]{in dieser} nur mit wenigen
Zeilen anzudeuten haben werde. Allein es ist von jeher mein gewissenhaft
befolgter Grundsatz gewesen, solche Andeutungen, die aufmerksame Leser in
jeder meiner Schriften in grosser Menge finden (sehen Sie z. B. meine \textit{Disquis.
Arithmet.} pag. 593[\textsuperscript{*})])\begingroup\renewcommand{\thefootnote}{}
\footnotetext{[\textsuperscript{*}) Werke I, S. 412, 413; es handelt sich dort um die Teilung des Lemniskate.]}
\endgroup\ stets dann erst zu machen, wenn ich den Gegenstand
für mich selbst ganz abgemacht habe, und so werden Sie übersehen, dass der
oben erwähnte Fall öfters vorkommen kann, wo um mit gutem Gewissen Eine
Zeile schreiben zu können, eine Monate erfordernde Meditation erfordert wird.
Diese Art zu arbeiten kann zuweilen die Folge haben, und hat sie zuweilen
gehabt, dass auf Dinge, die ich schon seit vielen Jahren besessen habe, später
ihrerseits auch andere kommen, und in der Bekanntmachung mir zuvorkommen;
sie wird vielleicht auch die Folge haben können, dass manches einmal mit mir
ganz untergeht, und ich weiss, dass einige meiner Freunde wünschen, dass ich
weniger in diesem Geiste arbeiten möchte: das wird aber nie geschehen; ich
kann einmal an lückenhaftem keine rechte Freude haben, und eine Arbeit,
an der ich keine Freude habe, ist mir nur eine Qual. Möge auch jeder in
\textls[120]{dem} Geiste arbeiten, der ihm am meisten zusagt. Was übrigens den gegenwärtigen
speciellen Fall betrifft, so hoffe ich, dass sich bei der Ausarbeitung
nicht so sehr viel finden wird, was noch lange aufhält.

Nach Vollendung dieser ersten Arbeit bin ich nicht abgeneigt, an ein
ausführlicheres Werk in deutscher Sprache über alle mit meinen Apparaten
anzustellenden Beobachtungen zu denken, worin denn auch diese Apparate
selbst vollständig beschrieben werden würden, was in der That allein schon
ein kleines Werkchen nöthig macht und von jener Abhandlung ausgeschlossen
bleiben muss, zumahl da es ohne Zeichnungen gar nicht geschehen könnte.
Alle Beobachtungen dieser Art gehören übrigens zu den reizendsten, die ich
kenne, sie übertreffen in dieser Beziehung auch die astronomischen, denen sie
an Praecision fast gleichkommen, ja in einiger Rücksicht noch übertreffen. Bei
der Winkelmessung kann man \(2''\) entschieden sichtbar machen, und es ist
\clearpage\SourceLeaf{90a}\HistoricalHead{85}{R}{GAUSS AN ENCKE.}
\noindent hauptsächlich nur der schwer ganz zu vermeidende Luftzug, welcher hindert,
dass man von dieser Genauigkeit entfernt bleibt, was indessen durch Vervielfältigung
wieder ersetzt werden kann. Bei allen Zeitansetzungen in Beziehung
auf Schwingungen hingegen ist die Schärfe entschieden weit grösser, als an
Beobachtungen am Pass[age] Instr[ument]. Es handelt sich immer nur von
Hunderttheilen der Secunde, nie, bei einiger Einübung, von Zehntheilen. Ich
würde selbst diese Schärfe für unglaublich halten, wenn ich sie nicht seit Monaten,
täglich vor mir sähe. Ich glaube, Sie würden es nicht bereuen, um
dies selbst zu sehen, eine Reise nach Göttingen gemacht zu haben, was man
ja jetzt in 40 Stunden kann. Wie sehr Sie mich durch einen solchen Besuch
erfreuen würden, brauche ich Ihnen nicht zu sagen. Mit dem Wunsche, bald
wieder durch einige Zeilen von Ihnen erfreut zu werden stets
\begin{flushright}
\begin{tabular}{c}
Ihr freundschaftlich ergebenster\\[.8ex]
C. F. \textsc{Gauss}.
\end{tabular}
\end{flushright}
\clearpage\endgroup
